% Document/LinearAlgebra/VectorSpaces/HomAndDual/HomSpaces.tex
% Phase 14: Dimension of Hom Spaces
% NOTE: Basic Hom vector space structure and Endomorphism Algebra moved to
%       Document/LinearAlgebra/VectorSpaces/LinearMaps/HomSpaces.tex (Chapter 2)
\section{Dimension of Hom Spaces}

\begin{proposition}[{$\Hom[K]{K}{V} \cong V$}]
\label{prop:hom-k-v}
    For any $K$-vector space $V$, $\Hom[K]{K}{V} \cong V$ via:
    \[
        \Phi: \Hom[K]{K}{V} \to V, \quad \Phi(T) = T(1_K)
    \]
    with inverse $\Phi^{-1}(v) = (a \mapsto av)$.
\end{proposition}

\begin{proof}
    \textit{$\Phi$ is well-defined}: For $T \in \Hom[K]{K}{V}$, $T(1_K) \in V$.
    
    \textit{$\Phi$ is linear}: 
    $\Phi(aS + bT) = (aS + bT)(1_K) = aS(1_K) + bT(1_K) = a\Phi(S) + b\Phi(T)$.
    
    \textit{$\Phi^{-1}$ is well-defined}: For $v \in V$, define $T_v: K \to V$ by $T_v(a) = av$.
    Then $T_v(ab + cd) = (ab + cd)v = a(bv) + c(dv) = aT_v(b) + cT_v(d)$, so $T_v \in \Hom[K]{K}{V}$.
    
    \textit{Inverse}: $\Phi(\Phi^{-1}(v)) = \Phi(T_v) = T_v(1_K) = 1_K \cdot v = v$.
    $\Phi^{-1}(\Phi(T)) = \Phi^{-1}(T(1_K)) = T_{T(1_K)}$.
    For any $a \in K$: $T_{T(1_K)}(a) = a \cdot T(1_K) = T(a \cdot 1_K) = T(a)$. Thus $T_{T(1_K)} = T$.
\end{proof}

\begin{remark}[{Asymmetry: $\Hom[K]{V}{K}$ is different}]
    Unlike $\Hom[K]{K}{V} \cong V$ which holds for all $V$, the space $\Hom[K]{V}{K}$ is
    \textbf{not} isomorphic to $V$ in general:
    \begin{itemize}
        \item For finite-dimensional $V$: $\Dim[K]{\Hom[K]{V}{K}} = \Dim[K]{V}$, but the isomorphism 
              depends on a choice of basis (non-canonical).
        \item For infinite-dimensional $V$: $\Dim[K]{\Hom[K]{V}{K}} > \Dim[K]{V}$, so no isomorphism exists.
    \end{itemize}
\end{remark}

\begin{corollary}[{$\Hom[K]{K}{K} \cong K$}]
\label{cor:hom-k-k}
    $\Hom[K]{K}{K} \cong K$.
\end{corollary}

\begin{proof}
    Apply Proposition~\ref{prop:hom-k-v} with $V = K$.
\end{proof}

\begin{theorem}[{Isomorphism for $\Hom[K]{U}{V}$}]
\label{thm:hom-dim-categorical}
    Let $U, V$ be $K$-vector spaces with bases $B_U$ and $B_V$ respectively. 
    
    If $\Dim[K]{V}$ is finite, then:
    \begin{align*}
        \Hom[K]{U}{V} &\cong \Hom[K]{K^{(B_U)}}{V} && \text{($U \cong K^{(B_U)}$)} \\
        &\cong \DirectProd{b \in B_U}{\Hom[K]{K}{V}} && \text{(contravariance: coproduct $\to$ product)} \\
        &\cong \DirectProd{b \in B_U}{\Hom[K]{K}{K^{(B_V)}}} && \text{($V \cong K^{(B_V)}$)} \\
        &\cong \DirectProd{b \in B_U}{\ExtDirectSum{b' \in B_V}{\Hom[K]{K}{K}}} && \text{(covariance, $\Dim[K]{V}$ finite)} \\
        &\cong \DirectProd{b \in B_U}{\ExtDirectSum{b' \in B_V}{K}} && \text{(Corollary~\ref{cor:hom-k-k})} \\
        &\cong \DirectProd{b \in B_U}{K^{(\Dim[K]{V})}} && \text{($\ExtDirectSum{b' \in B_V}{K} \cong K^{(\Card{B_V})}$)}
    \end{align*}
    Thus $\Hom[K]{U}{V} \cong \DirectProd{b \in B_U}{K^{(\Dim[K]{V})}}$, which has dimension 
    $\Card{K^{(\Dim[K]{V})}}^{\Dim[K]{U}}$.
\end{theorem}

\begin{corollary}[Finite-Dimensional Case]
\label{cor:hom-dim-finite}
    If $\Dim[K]{U} = m \in \bb{N}$ and $\Dim[K]{V} = n \in \bb{N}$, then:
    \[
        \Dim[K]{\Hom[K]{U}{V}} = m \cdot n
    \]
\end{corollary}

\begin{proof}
    By Theorem~\ref{thm:hom-dim-categorical}, $\Hom[K]{U}{V} \cong \DirectProd{i \in m}{K^{(n)}}$.
    Since both $m$ and $n$ are finite, finite products and coproducts coincide (biproduct), so 
    $\DirectProd{i \in m}{K^{(n)}} \cong K^{m \times n}$ and thus $\Dim[K]{\Hom[K]{U}{V}} = m \cdot n$.
\end{proof}

\begin{example}[{Dimension of $\Hom[K]{K^2}{K^3}$}]
    $\Dim[K]{\Hom[K]{K^2}{K^3}} = 2 \cdot 3 = 6$.
    
    This can be seen directly: a linear map $K^2 \to K^3$ is determined by where 
    it sends the basis vectors $(1,0)$ and $(0,1)$. Each can go to any vector in $K^3$,
    giving $3 + 3 = 6$ free parameters.
\end{example}

\begin{corollary}[Dimension of End]
\label{cor:end-dim}
    If $\Dim[K]{V} = n \in \bb{N}$, then $\Dim[K]{\End[K]{V}} = n^2$.
\end{corollary}

\begin{proof}
    $\Dim[K]{\End[K]{V}} = \Dim[K]{\Hom[K]{V}{V}} = n \cdot n = n^2$.
\end{proof}
